\section{Configurazione Iniziale}

\subsection{Step 1: Scarica i Modelli Ollama}

Prima di avviare Docker, scarica i modelli LLM necessari. Apri il terminale PowerShell:

\begin{lstlisting}
ollama pull qwen2.5:7b
ollama pull qwen2.5vl:3b
\end{lstlisting}

\textbf{Nota}: Se non hai Ollama installato localmente, puoi saltare questo step e i modelli verranno scaricati automaticamente al primo avvio del container (più lento).

\subsection{Step 2: Crea il File di Configurazione}

Nella radice della directory del progetto, crea un file chiamato \texttt{.env} con il seguente contenuto:

\begin{lstlisting}
NEO4J_USER=neo4j
NEO4J_PASSWORD=mirofish
LLM_MODEL=qwen2.5:7b
VISION_MODEL=qwen2.5vl:3b
MIROFISH_START_LOCAL_SERVICES=0
\end{lstlisting}

Questo file definisce:
\begin{description}
    \item[NEO4J\_USER] Nome utente per il database Neo4j
    \item[NEO4J\_PASSWORD] Password per accedere a Neo4j
    \item[LLM\_MODEL] Modello principale LLM
    \item[VISION\_MODEL] Modello per visione artificiale
    \item[MIROFISH\_START\_LOCAL\_SERVICES] Disabilita servizi locali (Docker li gestisce)
\end{description}
